\subsection{Probability measures and weak convergence}\label{subsec:preliminaries-probability}
\paraid{p-553202822150}
We collect the measure-theoretic tools used for averaging probabilities
and comparing distance operators on varying compact spaces.

\paraid{p-a5436302dabd}
For a metrizable space
$\ProbabilityCarrier $,
let
$\BorelSets(\ProbabilityCarrier )$
be its Borel
$\SigmaAlgebraPrefix$-algebra, namely, the
$\SigmaAlgebraPrefix$-algebra generated by the open subsets of
$\ProbabilityCarrier$.
A finite Borel measure
$\FirstMeasure$
on
$\ProbabilityCarrier$
is called \emph{Radon} if it is inner regular by compact sets,
meaning that for every
$\BorelSubset\in\BorelSets(\ProbabilityCarrier)$
we have
\[
 \FirstMeasure(\BorelSubset)
 =\sup\{\FirstMeasure(\CompactSubset)\mid
          \CompactSubset\subset\BorelSubset\text{ is compact}\}.
\]
We denote by
$\ProbabilityMeasures(\ProbabilityCarrier)$
the set of Radon measures with total mass one.
A \emph{probability law} on
$\ProbabilityCarrier$
is an element of
$\ProbabilityMeasures(\ProbabilityCarrier)$.
For Radon probabilities
$\FirstMeasure\in\ProbabilityMeasures(\ProbabilityCarrier)$
and
$\SecondMeasure\in\ProbabilityMeasures(\SecondProbabilityCarrier)$,
the notation
$\FirstMeasure\otimes\SecondMeasure$
denotes their product probability on
$\ProbabilityCarrier\times\SecondProbabilityCarrier$.
For every
$\BorelSubset\in\BorelSets(\ProbabilityCarrier)$
and
$\ClosedTestSet\in\BorelSets(\SecondProbabilityCarrier)$,
it satisfies
\[
 (\FirstMeasure\otimes\SecondMeasure)(\BorelSubset\times\ClosedTestSet)
 =\FirstMeasure(\BorelSubset)\SecondMeasure(\ClosedTestSet).
\]
These identities characterize the product probability.
For a Radon probability
$\FirstMeasure\in\ProbabilityMeasures(\ProbabilityCarrier)$,
its \emph{support} is
\[
 \supp(\FirstMeasure)
 =\{\BasePoint\in\ProbabilityCarrier\mid
    \FirstMeasure(\OpenSubset)>0\text{ for every open }\OpenSubset\subset\ProbabilityCarrier
    \text{ with }\BasePoint\in\OpenSubset\}.
\]
We say that
$\FirstMeasure$
has \emph{full support} if
$\supp(\FirstMeasure)=\ProbabilityCarrier$.
Equivalently,
\[
 \forall\OpenSubset\subset\ProbabilityCarrier\text{ open and nonempty}\quad
 \FirstMeasure(\OpenSubset)>0.
\]

\paraid{p-2c1ffb2b3240}
If
$\Mapping \colon \ProbabilityCarrier \to \SecondProbabilityCarrier $
is continuous and
$\FirstMeasure \in\ProbabilityMeasures(\ProbabilityCarrier )$,
its pushforward is the probability
$\Mapping \Pushforward\FirstMeasure $
defined for every
$\BorelSubset\in\BorelSets(\SecondProbabilityCarrier)$
by
\[
 (\Mapping\Pushforward\FirstMeasure)(\BorelSubset)
 =\FirstMeasure(\Mapping^{-1}(\BorelSubset)).
\]
For a nonempty metrizable space
$\ProbabilityCarrier$,
we write
\[
 \BoundedContinuousFunctions(\ProbabilityCarrier)
 =\{\TestFunction\colon\ProbabilityCarrier\to\RealNumbers\mid
       \TestFunction\text{ is continuous and bounded}\},
\]
and
\[
 \norm{\TestFunction}_\infty=\sup_{\BasePoint\in\ProbabilityCarrier}|\TestFunction(\BasePoint)|.
\]
For any compact metric space
$\BaseCarrier$,
we define
\[
 \ContinuousFunctions(\BaseCarrier)=\ContinuousFunctions(\BaseCarrier,\RealNumbers)
 =\{\TestFunction\colon\BaseCarrier\to\RealNumbers\mid\TestFunction\text{ is continuous}\},
\]
and
\[
 \norm{\TestFunction}_\infty=\max_{\BasePoint\in\BaseCarrier}|\TestFunction(\BasePoint)|.
\]
This is a real Banach space.
With continuous functions taking values in
$\ComplexNumbers$,
we obtain the complex Banach space
$\ContinuousFunctions(\BaseCarrier,\ComplexNumbers)$.

\paraid{p-bca1c91c25a9}
We equip
$\ProbabilityMeasures(\ProbabilityCarrier)$
with the weak topology,
the coarsest topology for which,
for every
$\TestFunction\in\BoundedContinuousFunctions(\ProbabilityCarrier)$,
the map
\[
 \ProbabilityMeasures(\ProbabilityCarrier)\to\RealNumbers
 \]
 by
 \[\FirstMeasure\longmapsto\int_\ProbabilityCarrier\TestFunction\,
 \IntegrationDifferential\FirstMeasure
\]
is continuous.
We write
$\FirstMeasure_\SequenceIndex\to\FirstMeasure$
weakly when for every
$\TestFunction\in\BoundedContinuousFunctions(\ProbabilityCarrier)$
we have
\[
 \int_\ProbabilityCarrier\TestFunction\,\IntegrationDifferential\FirstMeasure_\SequenceIndex
 \longrightarrow
 \int_\ProbabilityCarrier\TestFunction\,\IntegrationDifferential\FirstMeasure.
\]
On a compact space all continuous functions are bounded.
For bounded complex-valued continuous test functions the same convergence follows by
applying the definition to their real and imaginary parts.

\paraid{p-f85199b8afe5}
We construct averages of measures by specifying their integrals of continuous functions.
The Riesz--Markov--Kakutani theorem states that every positive normalized
linear functional on the real continuous functions on a nonempty compact
metric space is integration against a unique Radon probability measure.

\begin{theorem}[{Riesz--Markov--Kakutani, \cite[Theorem~7.10.4 and Corollary~7.10.5]{Bogachev2007II}}]\label{thm:riesz-markov-kakutani}
\paraid{p-37b86f9f43f7}
Let
$\MetricSpace{\BaseCarrier}{\MetricSymbol}$
be a nonempty compact metric space.
Let
\[
 \RepresentingFunctional\colon\ContinuousFunctions(\BaseCarrier,\RealNumbers)\to\RealNumbers
\]
be a linear functional.
Assume that for every
$\TestFunction\in\ContinuousFunctions(\BaseCarrier,\RealNumbers)$
with
$\TestFunction\geq0$
we have
\[
 \RepresentingFunctional(\TestFunction)\geq0.
\]
Let
$1$
denote the constant function with value one and assume that
\[
 \RepresentingFunctional(1)=1.
\]
Then there is a unique Radon probability
$\FirstMeasure\in\ProbabilityMeasures(\BaseCarrier)$
such that for every
$\TestFunction\in\ContinuousFunctions(\BaseCarrier,\RealNumbers)$,
we have
\[
 \RepresentingFunctional(\TestFunction)
 =\int_\BaseCarrier\TestFunction\,\IntegrationDifferential\FirstMeasure.
\]
\end{theorem}

\paraid{p-d1b3a84b65e1}
Positivity and normalization imply
$|\RepresentingFunctional(\TestFunction)|\leq\norm{\TestFunction}_\infty$,
so the functional
$\RepresentingFunctional$
is automatically continuous.
The cited representation theorem applies,
and evaluating at the constant function
$1$
proves that the representing measure
$\FirstMeasure$
has total mass one.

\begin{theorem}[{Portmanteau, \cite[Theorem~2.1]{Billingsley1999}, \cite[Theorem~8.2.3, Corollary~8.2.4(a)]{Bogachev2007II}}]\label{thm:portmanteau}
\paraid{p-e85c034ef572}
Let
$\ProbabilityCarrier$
be a metrizable space and let
$\FirstMeasure,\FirstMeasure_\SequenceIndex\in\ProbabilityMeasures(\ProbabilityCarrier)$
for
$\SequenceIndex\in\NonnegativeIntegers$.
Then the following conditions are equivalent.
\begin{enumerate}[label=\textup{(P\arabic*)},ref=\textup{(P\arabic*)}]
\item\label{item:portmanteau-weak}
\paraid{p-7c0d10bcd5a8}
For every
$\TestFunction\in\BoundedContinuousFunctions(\ProbabilityCarrier)$,
\[
 \lim_{\SequenceIndex\to\infty}\int_{\ProbabilityCarrier}\TestFunction\,\IntegrationDifferential\FirstMeasure_\SequenceIndex
 =\int_{\ProbabilityCarrier}\TestFunction\,\IntegrationDifferential\FirstMeasure.
\]
\item
\paraid{p-d9674e60f731}
For every closed subset
$\ClosedTestSet\subset\ProbabilityCarrier$,
\[
 \limsup_{\SequenceIndex\to\infty}\FirstMeasure_\SequenceIndex(\ClosedTestSet)
 \leq\FirstMeasure(\ClosedTestSet).
\]
\item
\paraid{p-533dc3ffdbf5}
For every open subset
$\OpenSubset\subset\ProbabilityCarrier$,
\[
 \FirstMeasure(\OpenSubset)
 \leq\liminf_{\SequenceIndex\to\infty}\FirstMeasure_\SequenceIndex(\OpenSubset).
\]
\end{enumerate}
\end{theorem}

\paraid{p-02818be37bab}
We also use the following consequence for measures on varying compact sets.

\begin{lemma}\label{lem:limit-support}
\paraid{p-a127ca5a2f91}
Let
$\FirstCompactSet_\SequenceIndex,\FirstCompactSet$
be nonempty compact subsets of a compact metric space
$\MetricSpace{\AmbientCarrier}{\AmbientMetric}$
with
$\HausdorffDistance{\AmbientMetric}(\FirstCompactSet_\SequenceIndex,\FirstCompactSet)\to0$.
If
$\FirstMeasure_\SequenceIndex\to\FirstMeasure$
weakly in
$\ProbabilityMeasures(\AmbientCarrier)$
and
$\FirstMeasure_\SequenceIndex(\FirstCompactSet_\SequenceIndex)=1$,
then
$\FirstMeasure(\FirstCompactSet)=1$.
\end{lemma}

\begin{proof}
\paraid{p-323518997d9a}
Define the function
$\TestFunction\colon\AmbientCarrier\to\RealNumbers$
by
\[
 \TestFunction(\IntegrationPoint)
 =\dist_{\AmbientMetric}(\IntegrationPoint,\FirstCompactSet).
\]
Since
$\AmbientCarrier$
is compact,
$\TestFunction$
is bounded and continuous.
Weak convergence and
$\FirstMeasure_\SequenceIndex(\FirstCompactSet_\SequenceIndex)=1$
imply
\[
 \begin{aligned}
 0\leq\int_\AmbientCarrier\TestFunction(\IntegrationPoint)
       \,\IntegrationDifferential\FirstMeasure(\IntegrationPoint)
 &=\lim_\SequenceIndex\int_\AmbientCarrier\TestFunction(\IntegrationPoint)
       \,\IntegrationDifferential\FirstMeasure_\SequenceIndex(\IntegrationPoint)\\
 &\leq\lim_\SequenceIndex
       \HausdorffDistance{\AmbientMetric}(\FirstCompactSet_\SequenceIndex,\FirstCompactSet)
 =0.
 \end{aligned}
\]
Since
$\TestFunction^{-1}(\{0\})=\FirstCompactSet$,
we obtain
$\FirstMeasure(\FirstCompactSet)=1$.
\end{proof}

\paraid{p-241d95ad9f1f}
Let
$\MetricSpace{\AmbientCarrier}{\AmbientMetric}$
be a compact metric space.
For
$\BorelSubset\in\BorelSets(\AmbientCarrier)$
and
$\SmallError>0$,
put
\[
 \BorelSubset^{\SmallError}
 =\{\BasePoint\in\AmbientCarrier\mid
       \dist_{\AmbientMetric}(\BasePoint,\BorelSubset)<\SmallError\}.
\]
For probabilities
$\FirstMeasure,\SecondMeasure\in\ProbabilityMeasures(\AmbientCarrier)$,
we define the \emph{L\'evy--Prokhorov distance} by
\[
 \ProkhorovDistance{\AmbientMetric}(\FirstMeasure,\SecondMeasure)
 =\inf\left\{\SmallError>0\;\middle|\;
 \begin{gathered}
 \text{for every }\BorelSubset\in\BorelSets(\AmbientCarrier),\\
 \FirstMeasure(\BorelSubset)\leq\SecondMeasure(\BorelSubset^{\SmallError})+\SmallError,\\
 \SecondMeasure(\BorelSubset)\leq\FirstMeasure(\BorelSubset^{\SmallError})+\SmallError
 \end{gathered}\right\}.
\]
This is a metric on
$\ProbabilityMeasures(\AmbientCarrier)$
and induces weak convergence
\cite[Section~6, pp.~72--73, Theorem~6.8]{Billingsley1999},
\cite[author's manuscript, Definition~1.12, Proposition~1.13 and Lemma~1.15]{ShioyaMetricMeasure}.

\begin{lemma}[{\cite[arXiv v5, Example~2.1(iii)]{Khezeli2023Framework}}]\label{lem:prokhorov-isometric-embedding}
\paraid{p-8ef64efd5d7f}
Let
$\MetricSpace{\BaseCarrier}{\MetricSymbol}$
and
$\MetricSpace{\AmbientCarrier}{\AmbientMetric}$
be compact metric spaces, and let
$\FirstEmbedding\colon\MetricSpace{\BaseCarrier}{\MetricSymbol}\to\MetricSpace{\AmbientCarrier}{\AmbientMetric}$
be an isometric embedding.
For all
$\FirstMeasure,\SecondMeasure\in\ProbabilityMeasures(\BaseCarrier)$,
we have
\[
 \ProkhorovDistance{\AmbientMetric}
 (\FirstEmbedding\Pushforward\FirstMeasure,\FirstEmbedding\Pushforward\SecondMeasure)
 =\ProkhorovDistance{\MetricSymbol}(\FirstMeasure,\SecondMeasure).
\]
\end{lemma}

% \paraid{p-d40306374206}

\paraid{p-b04d75056e40}
Every Borel probability on a Polish space is Radon
\cite[Theorem~7.1.7]{Bogachev2007II}.
We use the following theorem for Radon probability measures.

% \paraid{p-c3d38955d956}

\begin{theorem}[{\cite[Theorems~8.9.3 and~8.9.4]{Bogachev2007II}}]\label{thm:probability-topology}
\paraid{p-250872571fa5}
Let
$\ProbabilityCarrier$
be a separable metrizable space.
Then
$\ProbabilityMeasures(\ProbabilityCarrier)$
is separable and metrizable in the weak topology.
If
$\ProbabilityCarrier$
is Polish,
then
$\ProbabilityMeasures(\ProbabilityCarrier)$
is Polish.
If
$\ProbabilityCarrier$
is compact,
then
$\ProbabilityMeasures(\ProbabilityCarrier)$
is compact.
\end{theorem}

% \paraid{p-4be0d8c030c3}
% \paraid{p-489f0dc304a2}



\paraid{p-fd786a47d1fa}
Weak convergence is uniform over compact families of continuous test functions.

\begin{lemma}\label{lem:uniform-weak-integrals}
\paraid{p-97f9eb06751a}
Let
$\AmbientCarrier $
be a nonempty compact metric space,
let
$\FirstMeasure _\SequenceIndex \to\FirstMeasure $
weakly in
$\ProbabilityMeasures(\AmbientCarrier )$,
and let
$\CompactTestFamily \subset\ContinuousFunctions(\AmbientCarrier )$
be nonempty and compact in the uniform norm.
Then
\begin{equation}\label{eq:uniform-test-integrals}
 \sup_{\Mapping \in\CompactTestFamily }\left|\int_\AmbientCarrier  \Mapping \,\IntegrationDifferential \FirstMeasure _\SequenceIndex -\int_\AmbientCarrier  \Mapping \,\IntegrationDifferential \FirstMeasure \right|\to0.
\end{equation}
The same conclusion holds for complex continuous functions.
\end{lemma}

\begin{proof}
\paraid{p-3a52e680c1c2}
For each
$\SequenceIndex$,
choose
$\Mapping_\SequenceIndex\in\CompactTestFamily$
attaining the supremum in \eqref{eq:uniform-test-integrals}.
By compactness of
$\CompactTestFamily$,
every subsequence has a further subsequence
$\{\Mapping_{n_k}\}_{k\in\NonnegativeIntegers}$
converging uniformly to some
$\Mapping\in\CompactTestFamily$.
For these indices,
\[
 \left|\int_\AmbientCarrier\Mapping_{n_k}\,
           \IntegrationDifferential(\FirstMeasure_{n_k}-\FirstMeasure)\right|
 \leq2\norm{\Mapping_{n_k}-\Mapping}_\infty
      +\left|\int_\AmbientCarrier\Mapping\,
           \IntegrationDifferential(\FirstMeasure_{n_k}-\FirstMeasure)\right|
 \longrightarrow0.
\]
Thus the whole sequence of suprema tends to
$0$.
The same argument applies to complex functions,
since weak convergence controls both their real and imaginary parts.
\end{proof}

\paraid{p-9b10954ce9e8}
We use \Cref{lem:uniform-weak-integrals} to prove uniform convergence
of integrals whose integrands vary uniformly with compact parameters.

\begin{lemma}\label{lem:parameter-integrals}
\paraid{p-dc7ad8c030f6}
Let
$\CompactParameterSpace$
and
$\AmbientCarrier$
be nonempty compact metric spaces.
Let
$\FirstMeasure_\SequenceIndex\to\FirstMeasure$
weakly in
$\ProbabilityMeasures(\AmbientCarrier)$.
Let
$\Integrand_\SequenceIndex,\Integrand\colon\CompactParameterSpace\times\AmbientCarrier\to\ComplexNumbers$
be continuous real-valued or complex-valued functions.
Assume that
$\norm{\Integrand_\SequenceIndex-\Integrand}_\infty\to0$.
Then
\[
 \sup_{\ParameterPoint\in\CompactParameterSpace}
 \left|\int_\AmbientCarrier\Integrand_\SequenceIndex(\ParameterPoint,\IntegrationPoint)\,\IntegrationDifferential\FirstMeasure_\SequenceIndex(\IntegrationPoint)
 -\int_\AmbientCarrier\Integrand(\ParameterPoint,\IntegrationPoint)\,\IntegrationDifferential\FirstMeasure(\IntegrationPoint)\right|\to0.
\]
\end{lemma}

\begin{proof}
\paraid{p-03fa713f9e12}
The triangle inequality yields
\begin{equation}\label{eq:parameter-integral-error}
\begin{aligned}
&\sup_{\ParameterPoint\in\CompactParameterSpace}
 \left|
 \int_\AmbientCarrier
 \Integrand_\SequenceIndex(\ParameterPoint,\IntegrationPoint)
 \,\IntegrationDifferential\FirstMeasure_\SequenceIndex(\IntegrationPoint)
 -\int_\AmbientCarrier
 \Integrand(\ParameterPoint,\IntegrationPoint)
 \,\IntegrationDifferential\FirstMeasure(\IntegrationPoint)
 \right|\\
&\leq
 \norm{\Integrand_\SequenceIndex-\Integrand}_\infty
 +\sup_{\ParameterPoint\in\CompactParameterSpace}
 \left|
 \int_\AmbientCarrier
 \Integrand(\ParameterPoint,\IntegrationPoint)
 \,\IntegrationDifferential\FirstMeasure_\SequenceIndex(\IntegrationPoint)
 -\int_\AmbientCarrier
 \Integrand(\ParameterPoint,\IntegrationPoint)
 \,\IntegrationDifferential\FirstMeasure(\IntegrationPoint)
 \right|.
\end{aligned}
\end{equation}
The uniform error
$\norm{\Integrand_\SequenceIndex-\Integrand}_\infty$
in \eqref{eq:parameter-integral-error} tends to
$0$
by assumption.
The family
$\{\Integrand(\ParameterPoint,\cdot)\mid\ParameterPoint\in\CompactParameterSpace\}$
is compact in the uniform norm,
so \Cref{lem:uniform-weak-integrals} proves that the supremum of the integral differences
on the right-hand side of \eqref{eq:parameter-integral-error} tends to
$0$.
By \eqref{eq:parameter-integral-error},
the parameter integrals converge uniformly on
$\CompactParameterSpace$.
\end{proof}



\paraid{p-c858cdb38f3d}
We use the following continuous functions to detect the support of a measure
and approximate a distance profile by an integral.

\begin{lemma}\label{lem:metric-bump}
\paraid{p-e99ec25288a5}
Let
$\MetricSpace{\BaseCarrier}{\MetricSymbol}$
be a nonempty compact metric space,
let
$\FirstMeasure\in\ProbabilityMeasures(\BaseCarrier)$,
and let
$\Radius>0$.
For
$\BasePoint\in\BaseCarrier$,
define
\[
 \BumpFunction{\BaseCarrier}{\MetricSymbol}{\BasePoint}{\Radius}
 \colon\BaseCarrier\to[0,1]
\]
by
\begin{equation}\label{eq:metric-bump-definition}
 \BumpFunction{\BaseCarrier}{\MetricSymbol}{\BasePoint}{\Radius}(\ComparisonPoint)
 =\max\{1-\MetricSymbol(\BasePoint,\ComparisonPoint)/\Radius,0\}.
\end{equation}
Then for all
$\BasePoint,\BasePoint_0,\ComparisonPoint,\ComparisonPoint_0\in\BaseCarrier$,
\begin{equation}\label{eq:metric-bump-lipschitz}
 \left|\BumpFunction{\BaseCarrier}{\MetricSymbol}{\BasePoint}{\Radius}(\ComparisonPoint)
 -\BumpFunction{\BaseCarrier}{\MetricSymbol}{\BasePoint_0}{\Radius}(\ComparisonPoint_0)\right|
 \leq\frac{\MetricSymbol(\BasePoint,\BasePoint_0)
              +\MetricSymbol(\ComparisonPoint,\ComparisonPoint_0)}{\Radius}.
\end{equation}
Moreover,
\begin{equation}\label{eq:metric-bump-integral}
 \frac12\FirstMeasure\bigl(\MetricBall(\BasePoint,\Radius/2;\MetricSymbol)\bigr)
 \leq\int_\BaseCarrier
       \BumpFunction{\BaseCarrier}{\MetricSymbol}{\BasePoint}{\Radius}
       \,\IntegrationDifferential\FirstMeasure
 \leq\FirstMeasure\bigl(\MetricBall(\BasePoint,\Radius;\MetricSymbol)\bigr).
\end{equation}
\end{lemma}

\begin{proof}
\paraid{p-886f82531c59}
The function
$\SecondParameter\mapsto\max\{1-\SecondParameter/\Radius,0\}$
is
$1/\Radius$-Lipschitz.
The triangle inequality for
$\MetricSymbol$
therefore proves \eqref{eq:metric-bump-lipschitz}.
The pointwise bounds
\[
 \tfrac12\Indicator_{\MetricBall(\BasePoint,\Radius/2;\MetricSymbol)}
 \leq\BumpFunction{\BaseCarrier}{\MetricSymbol}{\BasePoint}{\Radius}
 \leq\Indicator_{\MetricBall(\BasePoint,\Radius;\MetricSymbol)}
\]
imply \eqref{eq:metric-bump-integral} by integration.
\end{proof}

\begin{lemma}\label{lem:normalized-metric-bump}
\paraid{p-ac4da2b9d373}
Let
$\MetricSpace{\BaseCarrier}{\MetricSymbol}$
be a nonempty compact metric space,
let
$\FirstMeasure$
be a probability measure on
$\BaseCarrier$
with full support,
and let
$\Radius>0$.
For every
$\BasePoint\in\BaseCarrier$,
the function
\begin{equation}\label{eq:normalized-bump-definition}
 \NormalizedBump{\BaseCarrier}{\MetricSymbol}{\BasePoint}{\Radius}{\FirstMeasure}
 =\frac{\BumpFunction{\BaseCarrier}{\MetricSymbol}{\BasePoint}{\Radius/2}}
 {\displaystyle\int_\BaseCarrier
      \BumpFunction{\BaseCarrier}{\MetricSymbol}{\BasePoint}{\Radius/2}
      \,\IntegrationDifferential\FirstMeasure}
\end{equation}
is well-defined and continuous,
and satisfies
\begin{equation}\label{eq:normalized-bump-properties-1}
 \NormalizedBump{\BaseCarrier}{\MetricSymbol}{\BasePoint}{\Radius}{\FirstMeasure}\geq0,
\end{equation}
and
\begin{equation}\label{eq:normalized-bump-properties-2}
 \int_\BaseCarrier
 \NormalizedBump{\BaseCarrier}{\MetricSymbol}{\BasePoint}{\Radius}{\FirstMeasure}
 \,\IntegrationDifferential\FirstMeasure=1,
\end{equation}
and
\begin{equation}\label{eq:normalized-bump-properties-3}
 \supp\left(\NormalizedBump{\BaseCarrier}{\MetricSymbol}{\BasePoint}{\Radius}{\FirstMeasure}\right)
 \subset\MetricBall(\BasePoint,\Radius;\MetricSymbol).
\end{equation}
\end{lemma}

\begin{proof}
\paraid{p-84a66f6444a3}
Fix
$\BasePoint\in\BaseCarrier$.
Since
$\supp(\FirstMeasure)=\BaseCarrier$,
\eqref{eq:metric-bump-integral} with radius
$\Radius/2$
proves that the denominator in \eqref{eq:normalized-bump-definition} is positive.
Continuity and \eqref{eq:normalized-bump-properties-1}--\eqref{eq:normalized-bump-properties-2}
follow from \eqref{eq:metric-bump-lipschitz} and \eqref{eq:normalized-bump-definition}.
By \eqref{eq:metric-bump-definition},
\[
 \supp\left(\NormalizedBump{\BaseCarrier}{\MetricSymbol}{\BasePoint}{\Radius}{\FirstMeasure}\right)
 \subset\{\ComparisonPoint\in\BaseCarrier\mid
              \MetricSymbol(\BasePoint,\ComparisonPoint)\leq\Radius/2\}
 \subset\MetricBall(\BasePoint,\Radius;\MetricSymbol).
\]
This proves \eqref{eq:normalized-bump-properties-3}.
\end{proof}

\paraid{p-0c1ccd024c12}
The next lemma compares product measures when one factor is a point mass.
It will be used to average laws over a convergent family of compact spaces.

\begin{lemma}[{\cite[Chapter~III, Lemma~1.1]{Parthasarathy2005}}]\label{lem:dirac-product-convergence}
\paraid{p-45f1a6509071}
Let
$\CompactParameterSpace$
be a compact metric space and let
$\ProbabilityCarrier$
be a Polish space.
Let
$\ParameterPoint_\SequenceIndex\to\ParameterPoint$
in
$\CompactParameterSpace$
and
$\FirstMeasure_\SequenceIndex\to\FirstMeasure$
weakly in
$\ProbabilityMeasures(\ProbabilityCarrier)$.
For a point
$\ParameterPoint\in\CompactParameterSpace$,
write
$\DiracProbability_\ParameterPoint$
for the probability concentrated at
$\ParameterPoint$.
Then
\begin{equation}\label{eq:dirac-product-convergence}
 \DiracProbability_{\ParameterPoint_\SequenceIndex}\otimes\FirstMeasure_\SequenceIndex
 \longrightarrow\DiracProbability_\ParameterPoint\otimes\FirstMeasure
 \quad\text{weakly in }\ProbabilityMeasures(\CompactParameterSpace\times\ProbabilityCarrier).
\end{equation}
\end{lemma}

% \paraid{p-c62c3005a39f}



\paraid{p-2cd0585c4b58}
For comparisons with metric measure geometry,
we also recall the box distance.

\begin{definition}\label{def:box-distance}
\paraid{p-b4f453559431}
An \emph{mm-space} is a complete separable metric space
$(X,d_X)$
with a Borel probability measure
$\mu_X$.
Two mm-spaces are \emph{mm-isomorphic} if there is a measure-preserving
isometry between their supports.
Let
$\MMClassSpace$
denote the set of these classes.
Put
$\BoxParameterInterval=[0,1)$
and let
$\BoxLebesgueMeasure$
be Lebesgue measure
on this interval.
A \emph{parameter} of
$X$
is a Borel map
$\BoxFirstParameter\colon\BoxParameterInterval\to X$
with
$(\BoxFirstParameter)_*\BoxLebesgueMeasure=\mu_X$.
Every mm-space has a parameter.
The \emph{box distance}
$\BoxDistance(X,Y)$
is the infimum of all
$\ErrorTolerance>0$
for which there are parameters
$\BoxFirstParameter$
of
$X$
and
$\BoxSecondParameter$
of
$Y$
and a Borel set
$\BoxLargeSubset\subset\BoxParameterInterval$
satisfying
\[
 \BoxLebesgueMeasure(\BoxLargeSubset)\geq1-\ErrorTolerance,
\]
and such that for every
$s,t\in\BoxLargeSubset$
we have
\[
 \left|d_X(\BoxFirstParameter(s),\BoxFirstParameter(t))
       -d_Y(\BoxSecondParameter(s),\BoxSecondParameter(t))\right|
 \leq\ErrorTolerance.
\]
It is a metric on
$\MMClassSpace$.
The topology it induces is called the \emph{box topology}
\cite[arXiv v1, Definitions~2.1, 2.6, 2.7 and Theorem~2.8]{KazukawaNakajimaShioya2024}.
\end{definition}
